\documentclass[a4paper,12pt]{article}

\usepackage[T1]{fontenc}
\usepackage[italian]{babel}
\usepackage{graphicx}
\usepackage{geometry}
\usepackage{tabularx}
\usepackage[table]{xcolor}
\usepackage[most]{tcolorbox}
\usepackage{fancyhdr}
\usepackage[hidelinks]{hyperref}

\newcommand{\groupmail}{alittlebyte.swe@gmail.com}
\newcommand{\grouplogo}{../../../.github/site-src/assets/logo_alittlebyte.png}
\newcommand{\groupsymbol}{../../../.github/site-src/assets/solo_logo.png}

\geometry{
    left=2.5cm,
    right=2.5cm,
    top=2.5cm,
    bottom=2.5cm,
    headheight=331pt,
    includeheadfoot
}

\pagestyle{fancy}
\fancyhf{}

\renewcommand{\headrulewidth}{0pt}
\fancyhead{}

\renewcommand{\footrulewidth}{0.4pt}
\fancyfoot[L]{\includegraphics[height=0.6cm]{\groupsymbol}\hspace{0.45em}aLittleByte}
    \fancyfoot[C]{\thepage}
\fancyfoot[R]{Verbale 2026-08-04}

\fancypagestyle{firstpage}{
    \fancyhf{}
    \renewcommand{\headrulewidth}{0pt}
    \renewcommand{\footrulewidth}{0.4pt}
    \fancyhead[C]{%
        \makebox[\textwidth][c]{%
            \begin{tabular}{c}
                \includegraphics[height=10cm]{\grouplogo}\\[1ex]
                \small\texttt{\groupmail}
            \end{tabular}%
        }
    }
    \fancyfoot[L]{\includegraphics[height=0.6cm]{\groupsymbol}\hspace{0.45em}aLittleByte}
        \fancyfoot[C]{\thepage}
    \fancyfoot[R]{Verbale 2026-08-04}
}

\providecommand{\glslink}[2]{\href{https://alittlebyte-19.github.io/Documentazione/glossario.html\#gls-#1}{\underline{#2}\textsuperscript{\scriptsize G}}}

\begin{document}

\thispagestyle{firstpage}

\vspace*{0.5cm}

\begin{center}
    {\LARGE \textbf{Verbale Interno - 2026-08-04}}
\end{center}

\vspace{1.5cm}

\begin{center}
    \renewcommand{\arraystretch}{1.3}
    \begin{tcolorbox}[
        width=0.92\textwidth,
        colback=gray!6,
        colframe=gray!45,
        boxrule=0.5pt,
        arc=2mm,
        boxsep=0pt,
        left=8pt, right=8pt, top=8pt, bottom=8pt
    ]
        \begin{tabularx}{\linewidth}{>{\bfseries}p{3.5cm} X}
            Gruppo      & aLittleByte (gruppo 19) \\
            Redattore   & Eleonora Bellet \\
            Verificatori & Leonardo Soligo\\
            Destinatari & Prof. Tullio Vardanega, Prof. Riccardo Cardin \\
            Data        & 2026-08-04\\
        \end{tabularx}
    \end{tcolorbox}
\end{center}

\newpage

\newgeometry{left=2.5cm,right=2.5cm,top=2.5cm,bottom=2.5cm,includefoot}
\tableofcontents
\newpage
\section{Informazioni Generali}
\section*{Specifiche Documento}
\hrule
\vspace{1cm}

\begin{tabular}{>{\bfseries}p{5cm} | p{10cm}}
Luogo Riunione & \glslink{discord}{Discord} \\[6pt]

Orario (Inizio - Fine) & 14:00 -- 14:30 \\[6pt]

\glslink{redattore}{Redattore} &
  E. Bellet\\[6pt]

\glslink{amministratore}{Amministratore} &
A. Gastaldello\\[6pt]

\glslink{verificatore}{Verificatore} &
L. Soligo\\[6pt]

Partecipanti &
  A. Oloni\\ &
  L. Soligo\\&
  A. Gastaldello\\&
  E. Bellet\\&
  D. Belluz\\&
  A. Maule\\&
  V. Y. Kaneda Fernandes\\[6pt]

\end{tabular}


\section{Ordine del Giorno}
    \begin{itemize}
    \item Integrazione branch: develop/main, check, code coverage e pipeline merge
    \item Strategia su app/specifiche: quando scrivere cosa e dipendenze dallo sviluppo
\end{itemize}


\section{Attività svolte}
\subsection{integrazione branch: develop/main, check, code coverage e pipeline merge}
Si aggiorna lo stato delle modifiche: Leonardo ha messo tutto in un unico branch di Github tutte le modifiche eseguite durante il periodo,
 restano check falliti e problemi su code coverage, branch e pipeline. L’idea è sistemare i check e poi permettere il merge senza disastri, eventualmente chiudendo branch non necessari. Si discute anche della sequenza rispetto alle modifiche

\subsection{Strategia su app/specifiche: quando scrivere cosa e dipendenze dallo sviluppo}
 Si discute la relazione tra refactoring dell’architettura, sviluppo dell’applicativo e stesura specifiche tecniche. 
 Viene indicato che la specifica tecnica è ancora in fase di standby, per il semplice fatto che dobbiamo aver completato la maggior parte dell'implementazione dell'MVP prima di continuare la stesura. I test e la parte delle norme/test vengono preferibilmente fatti alla fine, 
 perché dipendono dalla chiusura dell’applicativo backend/frontend. Si riconosce che scrivere la specifica implica anche dover modificare codice man mano che emergono incongruenze.

\section{To-Do List}
\begin{table}[h]
    \centering
    \begin{tabular}{|p{4.5cm}|p{5cm}|l|}
    \hline
        \rowcolor{lightgray}\textit{\textbf{Persona Incaricata}}  & \textbf{Task Assegnata} &\textbf{Link \glslink{issue}{Issue}} \\
    \hline
        \textit{A. Oloni, D. Belluz, L. Soligo} & Controllo dell'\glslink{analisi-dei-requisiti-adr}{Analisi dei Requisiti} &  \href{https://github.com/aLittleByte-19/Documentazione/issues/174}{\#174}\\
    \hline
        \textit{A. Gastaldello, L. Soligo} & Aggiornamento delle \glslink{norme-di-progetto}{Norme di Progetto} & \href{https://github.com/aLittleByte-19/Documentazione/issues/173}{\#173}\\
    \hline
        \textit{E. Bellet, A. Maule} & Aggiornamento del \glslink{piano-di-progetto}{Piano di Progetto} & \href{https://github.com/aLittleByte-19/Documentazione/issues/172}{\#172}\\
    \hline
        \textit{E. Bellet, A. Maule} & Aggiornamento del \glslink{piano-di-qualifica}{Piano di Qualifica} & \href{https://github.com/aLittleByte-19/Documentazione/issues/171}{\#171}\\
    \hline
        \textit{E. Bellet, A. Maule, V. Y. Kaneda Fernandes} & Aggiornamento delle Specifiche Tecniche & \href{https://github.com/aLittleByte-19/Documentazione/issues/150}{\#150}\\
    \hline
        \textit{D. Belluz, V. Y. Kaneda Fernandes} & Continuamento sviluppo dell'\glslink{mvp-minimum-viable-product}{MVP} & \href{https://github.com/aLittleByte-19/Documentazione/issues/166}{\#166}\\
    \hline
        \textit{E. Bellet} & Stesura \glslink{verbale-interno}{Verbale Interno} 04/08/2026 & \href{https://github.com/aLittleByte-19/Documentazione/issues/163}{\#163}\\
    \hline
    \end{tabular}
    \label{tab:placeholder}
\end{table}

\end{document}
